\mxcorrections{6}\label{C06-CORR}
\textbf{Usage professeur.} Les réponses qui suivent concernent les tâches non résolues dans le cours. Faire d’abord expliciter la démarche ; fournir une aide graduée avant la solution. Les durées sont des estimations à étalonner en classe.
\subsection*{Diagnostic et remédiation}
\begin{enumerate}
\item $6$. \textit{Action :} Reprendre le taux $2a+h$.
\item $\R\setminus\{2\}$. \textit{Action :} Annuler le dénominateur.
\item Positif à l’extérieur de $\intff{-2}{1}$, négatif entre les racines, nul aux bornes. \textit{Action :} Construire deux lignes de signes.
\item $x^3-2x^2+x-2$. \textit{Action :} Distribuer chaque terme sans omission.
\item Ils sont égaux. \textit{Action :} Substituer -x entre parenthèses.
\item $y=2$. \textit{Action :} Reprendre l’équation tangente et simplifier.
\end{enumerate}
\subsection*{Activité de recherche}
$x^3$ est strictement croissante avec tangente horizontale en 0. Pour une fonction croissante dérivable, les taux de part et d’autre sont positifs ou nuls, donc la dérivée ne peut être négative.
\subsection*{Applications immédiates}
\begin{enumerate}
\item $5x^4$ sur $\R$ ; $-\dfrac{2}{x^3}$ pour $x\ne0$ ; $\dfrac{1}{2\sqrt x}$ pour $x>0$.
\item $2x(3x-2)+3(x^2+1)=9x^2-4x+3$ ; $-\dfrac{2x}{(x^2+1)^2}$, sur $\R$.
\item $f\prime=2x-4$, donc décroissance jusqu’à 2 puis croissance. Minimum $-3$ en 2 ; valeurs aux bornes 1 et $-2$, donc maximum 1 en 0.
\item Première paire, seconde impaire ; différence $x^2+1-2x=(x-1)^2\geq0$, égalité en 1.
\item $A(x)=x(6-x)$ sur $\intff{0}{6}$ ; $A\prime=6-2x$. Maximum 9 m$^2$ en 3.
\end{enumerate}
\subsection*{Exercices : démarche et résultats}
\paragraph{Exercice 6.1 -- Une collection de dérivées.} $12x^3-2$ ; $2x(x-1)+(x^2+2)=3x^2-2x+2$ ; $\dfrac{[2(x+3)-(2x+1)]}{(x+3)^2}=\dfrac{5}{(x+3)^2}$ pour $x\ne-3$.
\paragraph{Exercice 6.2 -- Un tableau complet.} $f\prime=3(x-1)(x-3)$ ; signes $+,-,+$. Valeurs $f(0)=1$, $f(1)=5$, $f(3)=1$, $f(4)=5$. Minimum 1 en 0 et 3 ; maximum 5 en 1 et 4. Les bornes comptent.
\paragraph{Exercice 6.3 -- Un quotient.} $f\prime =1-\dfrac{4}{x^2}=\dfrac{(x-2)(x+2)}{x^2}$. Signe négatif avant 2, positif après. Minimum 4 en 2 ; pour $x>0$, $x+\dfrac{4}{x}\geq 4$, égalité seulement en 2.
\paragraph{Exercice 6.4 -- Comparer deux courbes.} $f-g=x(x-1)(x+1)$. Signes : négatif sur $\intoo{-\infty}{-1}$, positif sur $\intoo{-1}{0}$, négatif sur $\intoo{0}{1}$, positif sur $\intoo{1}{+\infty}$. Intersection en $-1,0,1$. On peut faire le tableau des trois facteurs sans dérivation.
\paragraph{Exercice 6.5 -- Un emballage.} $V=x(12-2x)^2=4x(6-x)^2$. $V\prime=12(x-2)(x-6)$, positif sur $\intfo{0}{2}$, négatif sur $\intoo{2}{6}$. $V(0)=V(6)=0$, $V(2)=128$ cm$^3$. Maximum pour x=2 cm ; boîte $8\times8\times2$.
\paragraph{Exercice 6.6 -- Stationnaire sans extremum.} $f\prime=3x^2$, donc valeur 0 en 0. Dans tout voisinage de 0, des x négatifs donnent $f(x)<0$ et des x positifs $f(x)>0$ ; aucun extremum.
\paragraph{Exercice 6.7 -- Comparer des stratégies.} Carré : $(x-3)^2+3\geq3$. Dérivée : $2x-6$, minimum en 3, valeur 3. La première évite une étude de variation complète ; la seconde illustre le lien entre chapitres.
\paragraph{Exercice 6.8 -- Reprise différée.} Elle décroît sur $\intoo{-\infty}{0}$ et $\intoo{0}{+\infty}$. Elle ne décroît pas sur leur réunion : $-1<1$ mais $f(-1)=-1<1=f(1)$. Le domaine n’est pas un intervalle.
\paragraph{CQFD, exercice 12 p. 156, questions 1 et 2.} $f(x)-g(x)=\dfrac12x^2-1=\dfrac12(x-\sqrt2)(x+\sqrt2)$. La courbe de f est strictement au-dessus de celle de g pour $x<-\sqrt2$ ou $x>\sqrt2$, au-dessous entre ces deux nombres ; égalité aux abscisses $\pm\sqrt2$.
\paragraph{CQFD, exercice 16 p. 156, questions f et g.} Pour $f(x)=x^3$, on peut choisir $\intoo{-\infty}{0}$ : $f(x)<0$ et la fonction est strictement croissante. Pour $g(x)=x^2-4$, on peut choisir $\intfo{0}{2}$ : $g(x)<0$ et la fonction est strictement croissante sur cet intervalle. Des sous-intervalles non réduits à un point conviennent aussi.
\subsection*{QCM : réponses et raisons}
\textbf{QCM 1.}
\begin{enumerate}
\item B. Puissance entière.
\item A. Taux de l’inverse.
\item D. Règle du produit.
\item C. Signe de dérivée et variation.
\item B. Dénominateur non nul.
\end{enumerate}
\textbf{QCM 2.}
\begin{enumerate}
\item C. $f(-x)=-f(x)$.
\item A. Les autres conclusions demandent plus.
\item D. Formule et domaine.
\item B. Atteint en 1.
\item C. $A\prime(x)=6-2x$ s’annule en 3 en passant du positif au négatif.
\end{enumerate}
\subsection*{Vrai ou faux}
\begin{enumerate}
\item Faux : $f(x)=\dfrac{1}{x}>0$ sur $\intoo{0}{+\infty}$, mais $f\prime(x)<0$.
\item Vrai, avec les hypothèses précisées ; aux bornes, la conclusion n’est pas imposée.
\item Faux sans hypothèses de signe : sur $\intff{-2}{-1}$, $u(x)=v(x)=x$ croissent mais $x^2$ décroît.
\end{enumerate}
\paragraph{Différenciation.} Pour un blocage technique, proposer la première étape seulement (domaine, facteur ou univers de référence), puis retirer l’aide. Pour un élève autonome, demander une seconde méthode et le traitement d’un cas limite. À l’oral, autoriser une reformulation initiale dans la langue maîtrisée, puis exiger une phrase mathématique précise en français.
